\documentclass[11pt,a4paper]{article}
\usepackage[utf8]{inputenc}
\usepackage[T1]{fontenc}
\usepackage[lithuanian]{babel}
\usepackage{lmodern}
\usepackage{geometry}
\geometry{a4paper, margin=2.5cm}
\usepackage{amsmath, amssymb}
\usepackage{enumitem}

\begin{document}

\begin{center}
    \Large \textbf{Kontrolinis darbas: Daugianariai} \\
    \large \textit{Apolonijaus varianto sprendimai}
\end{center}
\vspace{0.5cm}

\begin{enumerate}[label=\textbf{\arabic*.}]
    \item \textbf{Nustatykite duotų reiškinių tipus (R - racionalusis, T - trupmeninis, S - sveikas, I - iracionalusis).}
    \begin{enumerate}[label=\alph*)]
        \item $\displaystyle \frac{x^4 - \sqrt{5}}{3}$ \quad \textbf{Atsakymas: R, S} (nėra kintamojo vardiklyje ar po šaknimi).
        \item $\displaystyle \frac{4}{x^3 + 2}$ \quad \textbf{Atsakymas: R, T} (kintamasis yra vardiklyje).
        \item $\displaystyle x^{-4} + \sqrt{2} x$ \quad \textbf{Atsakymas: R, T} (nes $x^{-4} = \frac{1}{x^4}$, kintamasis yra vardiklyje).
        \item $\displaystyle \sqrt{x^2 - 9}$ \quad \textbf{Atsakymas: I} (kintamasis yra po šaknimi).
        \item $\displaystyle 2x+3^{-2} $ \quad \textbf{Atsakymas: R, S} (nes $3^{-2} = \frac{1}{9}$, kintamojo vardiklyje nėra).
        \item $\displaystyle \sqrt{2}x-7 $ \quad \textbf{Atsakymas: R, S} (kintamasis nėra po šaknimi, po šaknimi tik skaičius 2).
    \end{enumerate}

    \vspace{0.5cm}
    \item \textbf{Suprastinkite reiškinį ir nurodykite, ar tai vienanaris.}
    \[ \left( -3x^3y^{-2} \right)^3 \cdot \left( 3x^{-4}y^3 \right)^{-2} \]
    \textbf{Sprendimas:}
    \[ = \left( (-3)^3 \cdot (x^3)^3 \cdot (y^{-2})^3 \right) \cdot \left( 3^{-2} \cdot (x^{-4})^{-2} \cdot (y^3)^{-2} \right) \]
    \[ = \left( -27x^9y^{-6} \right) \cdot \left( \frac{1}{9}x^8y^{-6} \right) = -27 \cdot \frac{1}{9} \cdot x^9 \cdot x^8 \cdot y^{-6} \cdot y^{-6} = -3x^{17}y^{-12} \]
    \textbf{Atsakymas:} $-3x^{17}y^{-12}$ (arba $\displaystyle -\frac{3x^{17}}{y^{12}}$). Tai \textbf{nėra vienanaris}, nes kintamojo $y$ laipsnio rodiklis yra neigiamas (tai yra trupmeninis reiškinys).

    \vspace{0.5cm}
    \item \textbf{Apskaičiuokite reiškinio reikšmę:} \\
    \begin{align*}
        \frac{57^3+43^3}
        {0,1^{-2}\left(57\cdot14+43^2\right)}
        &= \frac{57^3+43^3}
        {100\left(57\cdot14+43^2\right)} \\
        &= \frac{(57+43)\left(57^2-57\cdot43+43^2\right)}
        {100\left(57\cdot14+43^2\right)} \\
        &= \frac{(57+43)\left(57(57-43)+43^2\right)}
        {100\left(57\cdot14+43^2\right)} \\
        &= \frac{100\left(57\cdot14+43^2\right)}
        {100\left(57\cdot14+43^2\right)} = 1.
    \end{align*}

    \vspace{0.5cm}
    \item \textbf{Įrodykite formules:}
    \begin{enumerate}[label=\alph*)]
        \item $\displaystyle x^m \cdot y^m = (x \cdot y)^m $
        \\ \textbf{Įrodymas:} Remiantis laipsnio apibrėžimu, kai $m$ yra natūralusis skaičius: $x^m = \underbrace{x \cdot x \cdots x}_{m \text{ kartų}}$ ir $y^m = \underbrace{y \cdot y \cdots y}_{m \text{ kartų}}$.
        Sudauginę gauname: $x^m \cdot y^m = (\underbrace{x \cdot x \cdots x}_{m}) \cdot (\underbrace{y \cdot y \cdots y}_{m}) = \underbrace{(xy) \cdot (xy) \cdots (xy)}_{m \text{ kartų}} = (xy)^m$.

        \item $(x-y)^3 = x^3 - 3x^2y + 3xy^2 - y^3$
        \\ \textbf{Įrodymas:} $(x-y)^3 = (x-y)(x-y)^2 = (x-y)(x^2 - 2xy + y^2) = x(x^2 - 2xy + y^2) - y(x^2 - 2xy + y^2) = x^3 - 2x^2y + xy^2 - x^2y + 2xy^2 - y^3 = x^3 - 3x^2y + 3xy^2 - y^3$.
    \end{enumerate}

    \vspace{0.5cm}
    \item \textbf{Raskite didžiausią ir mažiausią reiškinio $-3x^2 + \frac{3}{4} x + 1$ reikšmes.}
    \\ \textbf{Sprendimas:}
    \[
    -3x^2 + \frac{3}{4}x + 1
    = -3\left(x^2 - \frac{1}{4}x\right) + 1
    = -3\left(x^2 - 2 \cdot x \cdot \frac{1}{8} + \frac{1}{64} - \frac{1}{64}\right) + 1
    \]
    \[
    = -3\left(\left(x-\frac{1}{8}\right)^2 - \frac{1}{64}\right) + 1
    = -3\left(x-\frac{1}{8}\right)^2 + \frac{3}{64} + 1
    = -3\left(x-\frac{1}{8}\right)^2 + \frac{67}{64}
    \]
    \textbf{Atsakymas:} Didžiausia reikšmė yra $\frac{67}{64}$, kai $x=\frac{1}{8}$. Mažiausios reikšmės nėra.

    \vspace{0.5cm}
    \item \textbf{Reiškinį užrašykite pavidalu $a(x+b)^2 + c$ ir apskaičiuokite $\sqrt{-2a-4b+c+6}$.}
    \\ \textbf{Sprendimas:} Pradinis reiškinys:
    \[ (x-3)^3 - (x-1)^2(x-3) \]
    \[ = (x^3-9x^2+27x-27) - (x^3-5x^2+7x-3) \]
    \[ = x^3-9x^2+27x-27-x^3+5x^2-7x+3 = -4x^2+20x-24 \]
    \textbf{Išskiriame pilną kvadratą:}
    \[ -4x^2+20x-24 = -4(x-2,5)^2+1 \]
    Sulyginę su $a(x+b)^2+c$, gauname: $a=-4$, $b=-2,5$, $c=1$.
    Apskaičiuojame ieškomos šaknies reikšmę:
    \[ \sqrt{-2(-4) - 4(-2,5) + 1 + 6} = \sqrt{8 + 10 + 1 + 6} = \sqrt{25} = 5 \]
    \textbf{Atsakymas:} 5.

    \vspace{0.5cm}
    \item \textbf{Raskite parametrus $a$, $b$, $c$, kad lygybė $(ax+b)(x^2+2x-3) + cx = 2x^3+5x^2-8x-3$ būtų teisinga.}
    \\ \textbf{Sprendimas:} Atskliaudžiame kairiąją lygybės pusę:
    \[ ax^3 + 2ax^2 - 3ax + bx^2 + 2bx - 3b + cx = ax^3 + (2a+b)x^2 + (-3a+2b+c)x - 3b \]
    Kad daugianariai būtų lygūs su visomis $x$ reikšmėmis, jų koeficientai prie atitinkamų laipsnių turi sutapti su dešinės pusės ($2x^3+5x^2-8x-3$) koeficientais:
    \begin{itemize}
        \item Prie $x^3$: \quad $a = 2$
        \item Laisvasis narys: \quad $-3b = -3 \implies b = 1$
        \item Prie $x^2$: \quad $2a + b = 5$. Patikriname: $2(2) + 1 = 5$ (sutampa).
        \item Prie $x$: \quad $-3a + 2b + c = -8$. Įstatome $a$ ir $b$: \\
        $-3(2) + 2(1) + c = -8 \implies -6 + 2 + c = -8 \implies c - 4 = -8 \implies c = -4$.
    \end{itemize}
    \textbf{Atsakymas:} $a = 2$, $b = 1$, $c = -4$.
\end{enumerate}

\end{document}